% !PRINTED 697
\weberpage{697}
%<<CONT
Grund dieser Bemerkung kann man, wenn sich $n : \mu_n$ einem endlichen Grenzwerthe $\gamma$ nähert, $\alpha$ und $\beta$ diesem Grenzwerthe beliebig nahe annehmen, und erhält so die etwas präcisere Fassung des Theorems:

3. \gesperrt{Sind die positiven Zahlen $\mu_1, \mu_2, \mu_3, \dots$ so beschaffen, dass}
\[
\mathop{\mathrm{Lim}}\limits_{n=\infty} \frac{n}{\mu_n} = \gamma
\]
\gesperrt{ein endlicher Grenzwerth ist, so ist die unendliche Reihe}
\[
S = \frac{1}{\mu_1^s} + \frac{1}{\mu_2^s} + \frac{1}{\mu_3^s} + \cdots
\]
\gesperrt{für jedes $s$, was grösser als 1 ist, convergent, und es ist}
\[
\mathop{\mathrm{Lim}}\limits_{s=1} (s-1)\,S = \gamma.
\]

Es sei jetzt irgend ein System $M$ von \gesperrt{unendlich vielen} positiven Zahlen $\mu_n$ gegeben, die wir so ordnen
\weq{5}{\mu_1 \leqq \mu_2 \leqq \mu_3 \leqq \dots \qquad (M),}
so dass in der Reihe (5) niemals ein kleineres Glied auf ein grösseres folgt, und es bedeute $T$ die Anzahl von Gliedern in $M$, die nicht grösser als eine beliebig gegebene positive Grösse $t$ sind. Es gilt dann folgender Satz:

4. \gesperrt{Wenn einer der beiden Grenzwerthe}
\weq{6}{\mathop{\mathrm{Lim}}\limits_{n=\infty} \frac{n}{\mu_n}, \quad \mathop{\mathrm{Lim}}\limits_{t=\infty} \frac{T}{t}}
\gesperrt{endlich ist, so hat der andere denselben endlichen Werth.}

Es sei zunächst
\weq{7}{\mathop{\mathrm{Lim}}\limits_{n=\infty} \frac{n}{\mu_n} = \gamma}
ein endlicher Grenzwerth. Dann werden die $\mu_n$ mit unendlich wachsendem $n$ nothwendig ins Unendliche wachsen müssen, und es giebt für jedes positive $t$ einen Werth $m$, so dass
\weq{8}{\mu_m \leqq t < \mu_{m+1};}
$m$ wächst zugleich mit $t$ ins Unendliche und es ist $T = m$ [wegen (5)]. Daher nach (8)
\weq{9}{\frac{m}{\mu_m} \geqq \frac{T}{t} > \frac{m}{\mu_{m+1}}.}
